\section{Discussion}
\label{sec:discussion}

\subsection{Interpreting the Results}

\para{Near-ceiling performance limits diversity benefits.}
Our most striking finding is that both \haiku and \sonnet individually achieve
93.3\% accuracy on \advbench, a benchmark explicitly designed to be adversarial.
This high baseline leaves only 1--2 questions as potential improvement targets,
making statistical significance unattainable with $n{=}30$.
Modern Claude models trained with RLHF have likely encountered the classical
cognitive trap problems that form the basis of our benchmark~\citep{li2025system2},
making our benchmark insufficiently novel to expose meaningful failure mode differences.
A key lesson is that adversarial benchmarks for 2025-era models must use \emph{novel}
problems not likely present in training data.

\para{Error correlation confirms theoretical predictions.}
Despite the ceiling effect, we do observe a clear and statistically significant
difference in error correlation: the \haiku--\sonnet cross-capability pair
achieves $r{=}0.464$ (diversity index 0.536), while within-capability pairs
reach $r{=}0.695$ (diversity index 0.305).
This confirms \Htwo and is consistent with \citet{wu2023heterogeneous}'s
theoretical analysis: different capability levels do produce different failure modes.
The difference in diversity indices (0.536 vs.\ 0.305) is substantial in magnitude,
even though it does not translate to accuracy differences on our near-saturated benchmark.
With a more difficult benchmark (say, 50\% average accuracy), this difference in
error correlation would likely produce measurable accuracy gains.

\para{Safety alignment as a cross-cutting shared bias.}
The most important finding of this study is not about capability diversity at all:
it is about shared safety alignment.
Q27 reveals that all five conditions share an identical failure mode
because all Claude models are trained with similar RLHF principles
that refuse to give medical recommendations without clinical context.
This shared bias is \emph{more} harmful than correlated random errors,
because no within-family ensemble can overcome it.
\citet{rosales2025diverse}'s argument that models need \emph{complementary failure modes}
applies here: within-family models share not just capabilities but safety alignment properties,
making their failures systematically correlated on safety-sensitive questions.
Cross-family diversity (Claude + GPT + LLaMA) would be needed to overcome this.

\para{Sycophantic revision in the debate protocol.}
Our \debate protocol demonstrates the ``agreeability'' problem identified
by \citet{du2023debate}: \haiku accepts \sonnet's critique on Q23
and revises its originally correct answer to an incorrect one.
This ``sycophantic revision'' is concerning because our debate structure was
designed to mitigate it---\haiku proposes first and revises last,
giving it agency over the final answer.
Despite this structural choice, \haiku accepted \sonnet's more authoritative
but wrong interpretation.
Future debate protocols should include mechanisms to validate critique quality
before accepting revisions, \eg, requiring the critique to be accompanied
by a specific logical refutation rather than a general disagreement.

\para{Heterogeneous voting via haiku majority.}
An interesting side effect of the 2H+1S composition is that the \haiku majority
in the \hetero condition correctly overruled \sonnet's minority interpretation on Q23.
This is the \emph{opposite} of the expected ``wisdom of stronger models'' effect:
the stronger model (sonnet) was wrong, and the weaker model's majority protected
the ensemble from this error.
This illustrates that majority voting in heterogeneous ensembles is not simply
a weighted combination favoring the stronger model---it is sensitive to the
specific composition and question difficulty.

\subsection{Limitations}

\para{Small benchmark ($n{=}30$).}
With only 1--2 failures per condition, all McNemar tests are severely underpowered.
A meaningful comparison would require at least 200 questions at an average
accuracy of 70--80\% to achieve power of 0.80 for a 5 percentage point effect.
Our results should be treated as exploratory findings, not confirmatory evidence.

\para{Within-family diversity only.}
We compare \haiku and \sonnet only, both trained by Anthropic.
The core hypothesis about architectural diversity across families
(Claude vs.\ GPT vs.\ LLaMA) remains untested. Cross-family diversity
would expose providers' different training objectives, data mixes,
and safety implementations---the conditions most likely to yield
truly complementary failure modes.

\para{Temperature zero reduces sampling diversity.}
Running all conditions at temperature 0.0 means homogeneous ensembles produce
identical or near-identical responses, making 3$\times$\haiku effectively
equivalent to 1$\times$\haiku in most cases.
A more faithful study would use temperature $\approx$0.7 to introduce
productive within-model sampling diversity and separate scale effects
from capability effects.

\para{Classical adversarial tasks may be saturated.}
Problems like the bat-and-ball puzzle and the base-rate medical test
have been extensively discussed in LLM training data and benchmarks.
Truly novel adversarial tasks---\eg, programmatically generated
Knight-Knave configurations~\citep{wu2025debate}, synthetic mathematical
traps with novel surface patterns, or counterfactual causal chains---would
better test whether capability-level diversity benefits emerge in harder settings.

\para{Evaluation of open-ended answers.}
Our content-based matching for procedural questions (Q17, Q18) may miss
correct answers phrased differently.
In particular, Q18's answer extraction failure in \haiku could be a false negative:
\haiku's prose does mention the heat-based insight, suggesting the correct reasoning
is present but not surfaced in the final answer extraction.

\subsection{Implications for System Design}

For practitioners building multi-agent LLM systems,
our findings suggest three design principles:

\begin{enumerate}[leftmargin=*,itemsep=2pt,topsep=2pt]
    \item \emph{Use cross-family diversity when diversity is the goal.}
          Within-family capability diversity provides lower error correlation
          but shares safety alignment biases.
          For safety-sensitive applications, including models from different
          providers with different alignment training is essential.
    \item \emph{Prefer cascade over debate for heterogeneous systems.}
          Our debate results confirm that critique-revision protocols
          introduce sycophancy risk.
          Cascade architectures (route easy questions to the small model,
          escalate hard ones to the large model)~\citep{survey2025ensemble}
          avoid this failure mode while preserving cost efficiency.
    \item \emph{Benchmark difficulty must match model capability.}
          Evaluating 2025-era LLMs on classical cognitive trap problems
          provides a near-ceiling baseline that obscures diversity effects.
          Practitioners should regularly update their internal benchmarks
          to track genuinely challenging problems for their deployed models.
\end{enumerate}
